\section{5. Justification of the Functional Architecture}\label{ddjf-5}

\subsection{5.1. Basis for Functional Decomposition}\label{ddjf-5-1}

The GSEG architecture was developed on the basis of GSEG TRS, its RTM (Annex A) and functions F-01–F-08 in Chapter 4 of GSEG DDF. The primary decomposition criterion is separation of the automatic GSN–C\&C–DPS–EWS warning chain from MCC supervision, DAS archiving and GIS infrastructure services. The chapter addresses requirements for observation coordination, use of two-LEO data, threat assessment and automatic warning without blocking by the operator or archive.

\subsection{5.2. Justification of the Critical Functional Chain}\label{ddjf-5-2}

\subsubsection{5.2.1. GSN Communications Handling and Satellite Selection by C\&C}\label{ddjf-5-2-1}

GSN handles the SSEG–GSEG communications boundary, while C\&C handles selection of the LEO pair and reserve candidate, preparation of the two initial requests and assessment of their execution. Separating the communications function enables the logical correctness of coordination and the effectiveness of communications paths to be checked independently. C\&C uses previously prepared MCC orbital products without waiting for a new full orbital analysis or operator approval. This is justified by the two-second GSEG-REQ-PERF-001 requirement, although its compliance requires a test.

\textbf{Justification of the LEO reserve.} Identifying a reserve candidate during the first selection eliminates the need to rank the entire constellation again after detecting a failure, without routinely increasing the number of commands and observations from two to three. This is more favorable in terms of time than selecting a spacecraft only after a product is lost, and less burdensome for SSEG than tasking three spacecraft in parallel. This benefit can be realized only when the reserve forms a valid stereoscopic pair with each primary spacecraft, is reachable and is activated sufficiently early. For example, loss of one product near the end of the 30-second observation phase leaves no time to repeat that phase in full: \textbf{2+1 selection improves the chance of recovering an event but does not guarantee 60 s}. Agreement with SSEG is needed on the reserve status, replacement-request semantics, its execution time and \textbf{recalculation or validation of the onboard trajectory for the new pair}. Otherwise, valid products from individual LEO spacecraft may form an inconsistent DPS input set.

SSEG TRS provides for a LEO–LEO optical network (SSEG-REQ-ISL-001), GEO–LEO relays (SSEG-REQ-ISL-002) and autonomous packet routing (SSEG-REQ-ISL-006). Functionally, this means that GSN does not need to maintain direct radio contact with both selected LEO spacecraft. A reachable satellite node connected by an operational route to the command recipients and the sources of the required products is sufficient. The preferred physical design uses two separate GSN paths to two visible GEO spacecraft, providing an alternative entry point to the ISL network. This is a P design solution, not the result of a completed visibility or latency analysis. Operation remains conditional on preservation of critical-message priority and compliance with SSEG-REQ-PERF-003/\allowbreak{}004 timing limits.

\subsubsection{5.2.2. Justification of Separating Threat Assessment (DPS) and Warning Dissemination (EWS)}\label{ddjf-5-2-2}

The division into DPS and EWS follows from the different correctness criteria for the two functions. \textbf{DPS is responsible for the substantive validity of a warning}, including use of the complete set of valid observations, uncertainty assessment and threat qualification (GSEG-REQ-FUN-003–005). \textbf{EWS is responsible for faithful and timely delivery of the result of that qualification} to the appropriate recipient (GSEG-REQ-IF-002, PERF-002). If dissemination independently interpreted observation data or recalculated the threat threshold, there would be two places where the same decision was made. Divergent rules could result in a required warning not being sent or a message being sent without justification in the DPS assessment. Separating responsibilities reduces this risk and allows changes to message delivery independently of assessment logic.

This choice is particularly important for an \textbf{ambiguous result} and \textbf{inability to assess}. The former may qualify for a warning despite uncertainty; the latter does not confirm safety but denotes insufficient data. EWS should therefore receive an unambiguous qualification outcome from DPS associated with the event identifier, uncertainty assessment and appropriate product version; it should not reconstruct the decision from incomplete message fields. This condition follows from FUN-003/\allowbreak{}005 and DATA-001/\allowbreak{}003, but the detailed exchange format remains subject to the ICD.

The alternative of a single undivided assessment and transmission function could eliminate one internal interface but would more tightly couple changes to the assessment algorithm, recipient system and failure handling. \textbf{The cost of the selected division} is an additional exchange boundary that must be protected against delay, loss, repetition or product-version inconsistency. It is therefore insufficient to state that both blocks exist: an integrated test must demonstrate a correct decision and delivery within the \textbf{15 s} limit, including during GIS domain failover and for all qualification outcomes. The functional separation is considered justified at this stage, but its timing effectiveness and resilience to incorrect decision transfer remain to be confirmed.

\subsection{5.3. Justification of Separating Supporting Functions}\label{ddjf-5-3}

\textbf{MCC — data preparation and supervision outside the first-decision path.} An alternative would be to calculate orbital geometry each time and wait for an operator decision after the GEO report. This would be difficult to reconcile with the \textbf{2 s} limit and would introduce a dependency of warning time on personnel and planning-tool availability. Therefore, advance preparation of orbital products by MCC and their automatic use by C\&C were selected; this is also consistent with GSEG-REQ-ARC-002. The cost of this solution is the risk of using an outdated prediction. Justification will be complete only after defining product validity and checking how C\&C responds to absence or expiry of the products. Separating MCC does not mean that prolonged loss of orbital updates is immaterial to the mission.

\textbf{DAS — preserving timeliness by separating two objectives.} Requirements for event history and retention (DATA-004/\allowbreak{}005) and for rapid first warning (PERF-002, ARC-002) have different timing characteristics. If completion of archiving were a condition for transmission, storage delay or failure could stop the warning despite a correct DPS assessment. Removing this condition from the critical path eliminates an unnecessary dependency and allows separate verification of warning time and archive completeness. It does not remove the obligation to preserve the history — both properties must be demonstrated in appropriate tests. The detailed recording mechanism remains in lower-level designs; no additional buffering technology is selected at this stage.

\textbf{GIS — shared infrastructure with protection of critical functions.} Completely separate infrastructure for each subsystem would mitigate some shared dependencies, but would duplicate computing, network and maintenance resources and make consistent configuration and time management more difficult. The common GIS platform is therefore a justified segment-level choice \textbf{provided that} MCC/\allowbreak{}DAS workloads do not deprive C\&C/\allowbreak{}DPS/\allowbreak{}EWS functions of resources and that failure of a single equipment domain does not stop the entire chain. This decision is supported by PERF-001/\allowbreak{}002 and REL-001/\allowbreak{}003; Chapter 6 analyzes its physical implementation and common-cause hazards.

In particular, priority allocation of a limited DPS resource pool is justified by more predictable behavior than uncontrolled concurrency or a simple FIFO queue. \textbf{It does not yet justify the number of 10 events} as sufficient for actual mission traffic. To confirm the choice, assessment-time measurements and saturation-scenario tests are required, including loss of one GIS domain; the option-comparison approach is given in 5.5 and the analysis plan in Chapter 7.

\subsection{5.4. Requirements Allocation}\label{ddjf-5-4}

The decomposition remains consistent with the GSEG TRS RTM: FUN-001/\allowbreak{}002 and PERF-001 — C\&C/\allowbreak{}GSN/\allowbreak{}GIS; FUN-003–006 and DATA-002/\allowbreak{}003 — DPS, together with EWS for FUN-005; FUN-007 — MCC/\allowbreak{}EWS; FUN-008 — DPS/\allowbreak{}EWS; FUN-009 — a constraint on the entire GSEG. ARC-001/\allowbreak{}002 cover location and chain independence, while IF-001–003 cover GSN, EWS and MCC/\allowbreak{}C\&C interfaces. Archiving requirements DATA-004–006 are allocated to DAS; SEC concerns the applicable blocks and cross-cutting protection; REL-001–003 cover the entire critical chain including GIS, local gateways and connections controlled by GSEG. OPS-001/\allowbreak{}002 concern MCC, while VER-001–003 concern integration and verification of integrated GSEG. Annex A RTM contains the detailed allocation of each requirement.

\subsection{5.5. Options and Decisions Requiring Justification}\label{ddjf-5-5}

For communications, (1) direct GSN contact with each selected LEO spacecraft and (2) routing via satellite nodes, in particular two GEO spacecraft acting as gateways to the ISL network, were considered. The second option is preferred because it reduces dependence on short direct-visibility windows of LEO spacecraft and exploits capabilities provided for in SSEG TRS. However, it does not automatically provide two independent routes within the ISL network itself and therefore requires analysis of the dynamic connectivity graph, terminal allocation, QoS queues and node failures.

\textbf{Overload-handling options:} unlimited parallel processing is not supported by fixed computing capacity, while a simple FIFO queue may block an event requiring urgent assessment behind less important events. Therefore, \textbf{a limited pool of 10 DPS slots and a priority-allocation queue} (P) were selected. C\&C assigns preliminary priority from the GEO report and event deadline, and DPS adjusts it on the basis of valid LEO data. Regional priority must not replace classification or authorize silently abandoning global events; queue-age rules and explicit saturation reporting are required. The architectural advantage is predictable resource use after loss of one domain; the cost is \textbf{the lack of evidence of compliance with 60 s after capacity is exceeded}. As the current TRS does not establish an upper event-arrival limit, such overload cannot be considered a permissible exception without a decision at system requirements level. Tests should compare immediate arrivals of 10, 11 and 20 events and the handling order of a high-priority event arriving later. The number 10 itself remains an unproven sizing hypothesis, not a mission limit.

\subsection{5.6. Justification Result}\label{ddjf-5-6}

The functional decomposition is consistent with TRS. The 2+1 option reduces the cost of selecting LEO spacecraft again, and the priority-based pool of 10 assessments prevents uncontrolled resource allocation, but neither benefit has yet been confirmed in terms of timing. Chapters 7 and 8 describe the evidence required for closure.

\par\noindent{\color{TableBorderBlue}\rule{\linewidth}{0.4pt}}\par
